\uuid{s8cR}
\exo7id{1565}
\titre{exo7 1565}
\auteur{legall}
\organisation{exo7}
\datecreate{1998-09-01}
\isIndication{false}
\isCorrection{false}
\chapitre{Endomorphisme particulier}
\sousChapitre{Endomorphisme orthogonal}
\cpge{mpsi-mp2i.calcul-matriciel-systemes.e, pcsi.determinants.c, ptsi.determinants.a, mpsi-mp2i.espaces-prehilbertiens.c, pcsi.espaces-prehilbertiens.c, bcpst1.matrices, bcpst1.geometrie.f, mp-mpi.endomorphismes-euclidien.b, pc.endomorphismes-euclidiens.b, psi.endomorphismes-euclidien.b, pt.isometries-euclidien.b, bcpst2.produit-scalaire.a}
\module{Algèbre}
\niveau{L2}
\difficulte{}

\contenu{
\question{
Soit $  A=(a_{i,j})\in \mathcal{O}(n,\R)  .$ Montrer pour tout $  j \in \{
1, \cdots , n\}  $ l'\'egalit\'e~: $  \displaystyle{ \sum _{i=1}^n a_{i,j}^2=1  .}$ En d\'eduire
que si $  A  $ est triangulaire sup\'erieure elle est diagonale.
}
}
